%%%%%%%%%%%%%%%%%%%%%%%%%%%%%%%%%%%%%%%%%%%%%%%%%%%%%%%%%%%%%%%%%%%%%%%%%%%%%%%%
% DISEÑO E IMPLEMENTACIÓN %

\cleardoublepage
\chapter{Diseño e implementación}
\label{sec:diseno}


\section{Fase 1: Definición del proyecto}
\label{sec:definicion-proyecto}

El punto de partida del proyecto fue una reunión con el tutor en abril de 2025, en la que le comenté posibles ideas de trabajo.
En esa primera toma de contacto surgieron un par de ideas más, las cuales finalmente se acabaron descartando en nuestra segunda reunión,
ya que decidimos adaptarnos un poco más a la situacion actual con el tema de las paginas webs y automatizaciones.


La primera idea consistía en una aplicación web que, a partir de datos de APIs hospitalarias, ayudara a los pacientes a encontrar el centro 
médico más adaptado a sus necesidades según una serie de parámetros. 
La segunda idea ampliaba esa primera hacia un asistente de mudanza. Imaginese que alguien quiere instalarse en una ciudad, la aplicación
recogería todos los datos de las APIs publicas y un parseo de los datos proporcionados por agencias y otras empresas, de tal forma que 
introduciendo el usuario sus requisitos, se le organizaría un posible plan de vida en la zona.
Ambas ideas tenían en común la lectura y procesamiento de APIs publicas, pero el alcance era demasiado limitado, lo que limitaba el valor del 
proyecto. 


En la segunda reunion, realizada en Julio de 2025, llegamos a una idea más acotada, una aplicacion web Django encargada de procesar los datos
de cualquier API publica y generar automaticamente un sitio web personalizado. Mas tarde pensé en posibles nombres para el proyecto, hasta que 
al final dimos con el nombre de WebBuilder.


En ese momento aún estaba liado con los examenes, pero en los ratos libres iba recogiendo informacion y creando las posibles fases del sistema:
validación de la URL, procesamiento de los datos recogidos, preferencias sobre el diseño, etc. David valoró positivamente estos pasos, lo que 
me dio mas energía para seguir adelante con el proyecto.


Durante los meses de verano el proyecto comenzo poco a poco creando las partes mas sencillas para ir dandole forma: la pagina de inicio, el 
sistema de login y registro de usuarios. Aunque en el futuro todo esto cambiaría, la base mas  o menos se mantendría.

% Poco a poco conforme el proyecto avanzo fui teniendo problemas para poder hacer de verdad  algo personalizado y diferente, era muy dificil procesar los datos de todas las APIs y que no saliesen datos genericos, ahi le surgió a David una idea, la cual dio el giro completo del proyecto y lo convirtio en algo innovador.

% Ibamos a incorporar un modulo de Inteligencia Artificial, la idea no solo era procesarle los datos y que el eligiese los mas adecuados, sino construir un sistema en el que la IA fuese el motor principal y actuase como un desarrollador guiado con prompting engenieering. Pero este cambio estara detallado mas adelante.

Durante los primeros meses siguientes no hubo un gran desarrollo de codigo, sino que dedicamos este periodo a definir con mas detalle el flujo
de la aplicacion, estuvimos investigando tecnologias y diseñando procesos que deberia seguir el sistema. Este trabajo previo fue muy importante
para el correcto desarrollo guiado de la aplicacion.


\section{Fase 2: Comienzo del desarrollo}
\label{sec:comienzo-desarrollo}

Con la propuesta del proyecto ya definida, el desarrollo comenzó con fuerza en febrero de 2026. Poco a poco fui creando la estructura base de 
la aplicación: modelos, sistema de ingestion de datos, logica de analisis heuristico, y una interfaz minima y bastante pobre para poder hacer
pruebas del flujo completo. La arquitectura hasta este punto era bastante sencilla: una aplicacion Django que recibia una URL, descargaba los 
datos de la API, los analizaba y devolvía una propuesta de estructura para el futuro sitio web.

\subsection{Modelo de datos}
\label{subsec:modelo-datos}

El unico modelo de la base de datos en esta fase era \texttt{APIRequest}, que recogia los datos de cada peticion URL realizada.
Se encargaba de recoger los siguientes campos: Datos en crudo recogidos de la API (\texttt{raw\_data}), la estructura de datos ya parseada (\texttt{parsed\_data}),
el estado del análisis mediante un campo de elección (\texttt{pending}, \texttt{processed} o \texttt{error}), y el mapping resultante que se mostraria al usuario
(\texttt{field\_mapping}). Este modelo fue el eje central durante esta fase y será la base del análisis durante todo el proyecto

\subsection{Ingestion de datos}
\label{subsec:ingestion-datos}

El modulo de ingestion de datos, ubicado en \texttt{utils/ingest/} es el encargado de validar la URL introducida por el usuario, realizar la descarga del contenido
de forma segura y detectar automaticamente el formato de los datos. La validacion comprobaba que la URL procederia de \texttt{http} o \texttt{https}, que fuese
procedente de un dominio valido y que su longitud no superase los 2048 caracteres. La descarga se realizaba en modo streaming, es decir, en vez de esperar que
el contenido del archivo se descargue entero, vamos leyendo trozo a trozo procesado, cada trozo tiene un limite de 1MB, y un timeout de 8 segundos. Conseguimos de 
esta forma evitar que una API con respuestas excesivamente grande nos tumbe o bloquee el sistema.

La deteccion de formato se realiza inspeccionando los primeros caracteres del contenido de la API, de esta forma conseguiamos distinguir entre JSON y XML, que son 
los contenidos manejados hasta este momento, mas adelante en el capitulo~\ref{subsec:nuevos-formatos-entrada} se aumenta el tipo de datos a procesar. En cualquier caso 
siempre se devuelve una estructura Python homogenea, que se utiliza en adelante para el sistema, de forma que no es necesario seguir trabajando con el formato 
original.

\subsection{Deteccion de la colección principal}
\label{subsection:deteccion-coleccion-principal}

La parte mas compleja de esta fase fue entender la estructura interna de los campos de las distintas APIs. Los campos de las APIs no suelen ser una lista plana de items
en la raiz, es decir, la colección principal no esta expuesta, sino que esta envuelta bajo claves de metadatos como \texttt{results}, \texttt{data} o \texttt{items},
o incluso con varios niveles de profundidad. 

Para resolver esto se implemento un nuevo modulo: \texttt{utils/analysis}, donde se crearon ficheros con funciones importantes como \texttt{find\_main\_items}, la cual recorria recursivamente la estructura de la API, y se evaluaba con un sistema de puntuacion propio para sacar la estructura principal.

Cada lista candidata se puntuaba segun cinco factores:

\begin{itemize}
    \item \textbf{Densidad de diccionarios}: porcentaje de elementos de la lista que son diccionarios. Las listas con menos de un 55\% de dicts recibían una penalización proporcional, ya que era poco probable que representaran una colección de entidades reales.
    \item \textbf{Tamaño}: listas con más items recibían mayor puntuación, aunque con un límite para evitar que una lista enorme pero trivial se impusiera sobre una colección más pequeña pero más relevante.
    \item \textbf{Consistencia de claves}: si todos los diccionarios de la lista compartían aproximadamente las mismas claves, era muy probable que representaran entidades del mismo tipo. Se calculaba como el promedio de claves propias de cada item sobre la unión de todas las claves observadas.
    \item \textbf{Bonus por claves semánticas}: si los campos de la colección incluían nombres típicos de contenido como \texttt{title}, \texttt{name}, \texttt{id},    \texttt{description} o \texttt{image}, se añadía una bonificación proporcional al número de coincidencias.
    \item \textbf{Penalización por profundidad}: los paths muy anidados o que pasaban por índices numéricos se penalizaban para favorecer colecciones cercanas a la raíz, que solían ser las más relevantes.
\end{itemize}

Al finalizar este scoring, se seleccionaba la lista con mayor puntuacion como coleccion principal, junto con la ruta dentro de la estructura original.

\subsection{Sugerencia automatica}
\label{subsec:sugerencia-automatica}

Una vez localizada la coleccion principal, se realizo la implementacion del nuevo fichero \textbf{utils/mapping.py}, encargado de sugerir automaticamente que campos
del dataset debian ocupar cada rol: fecha, id, descripcion, etc. Para ello se utilizaba un catalogo predefinido que asociaba cada rol con sus nombres más habituales.
Por ejemplo un campo llamado \texttt{published\_at}, obtenia un alto score para el rol fecha ya que contenia la subcadena \texttt{published}.
Ademas de la comparacion por nombre, el sistema añadia score extra si el tipo del campo observado coincidia con el tipo esperado en el rol.

Para evitar que la misma clave del dataset se asignara como primera sugerencia de varios roles a la vez, se implementó una funcion llamada \texttt{suggest\_mapping}, que aplicaba una asignación por orden de prioridad: primero se reservaba la mejor clave para el rol \texttt{id}, luego para \texttt{title}, luego para \texttt{description}, y así sucesivamente. Esto reducía la repetición de sugerencias y hacía el mapping resultante más útil para el usuario, sin embargo, ocurrian errores como que algunos campos se quedasen vacios o que al asignarse por orden de prioridad, campos sueltos se asignasen a roles
que no tenian relacion.

\subsection*{Eleccion pagina}
\label{subsec:eleccion-pagina}

Para guiar mejor al usuario durante el proceso de mapping, se añadió un sistema de perfiles de intención en \texttt{intents.py}. El usuario podía indicar si quería construir un blog, un portfolio, un catálogo o un directorio, y el sistema reordenaba los roles y ajustaba cuáles eran obligatorios, recomendados u opcionales según ese tipo de sitio.


\subsection*{Calidad del mapping}
\label{subsec:calidad-mapping}
Además, se implementó un sistema de puntuación de calidad del mapping en \texttt{mapping.py}, que evaluaba de 0 a 100 el mapping introducido por el usuario teniendo en cuenta si había asignado los roles más importantes, si existían duplicados entre roles críticos y qué campos quedaban vacíos. El resultado se mostraba al usuario con una etiqueta cualitativa: \emph{Excelente}, \emph{Bueno}, \emph{Aceptable} o \emph{Mejorable}.


\subsection*{Tests}
\label{subsec:test1}

Desde la primera fase se escribieron tests para los módulos principales: análisis, mapping, normalizer y parsers. Esto permitió verificar el 
comportamiento del sistema de scoring con datasets variados, detectando casos en los que las heurísticas fallaban antes de que llegaran al usuario.
\subsection*{Limitaciones del enfoque heurístico} \label{subsec:limitaciones1}

Conforme se probaba el sistema con datasets reales muy distintos entre sí, los límites del enfoque heurístico se fueron haciendo evidentes. Una API de plantas medicinales, una de personajes de ficción y una de productos de e-commerce tienen nomenclaturas de campos completamente diferentes, y un sistema de reglas fijas no podía cubrir todos los casos con la misma fiabilidad. Cuando los nombres de los campos no coincidían con los campos de la API, el mapping era genérico y malo para el usuario. El sistema funcionaba razonablemente bien con APIs convencionales y con campos estándar, pero fallaba ante cualquier dataset con nombres o estructura poco habituales. Esta limitación fue la que motivó el cambio de enfoque descrito en la siguiente fase.


\section{Fase 3: Integracion LLM y n8n}
\label{sec:integracion-llm-n8n}

Con la fase 2 terminada y viendo que las limitaciones eran evidentes, decidimos hacer un cambio muy importante para el futuro del proyecto. Abandonamos 
la idea del sistema de reglas que teniamos, aunque no en la totalidad, ya que había partes que nos seguirían siendo útiles. Decidimos así integrar un
modelo de lenguaje (LLM) como nucleo o motor del análisis y generación de las paginas. Hasta este momento el sistema intentaba entender los datos mediante
puntuaciones y comparaciones, pero decidimos que fuese la Inteligencia Artificial la encargada de escoger y proponer la estructura del sitio y decidiera que 
campos mostrar. Como comente antes, no todo el flujo de análisis e ingesta se cambio, se eliminaron las partes de score y propuesta de dataset, pero en parseo 
y análisis de los campos se mantuvo y modifico levemente, se metió todo el dataset en una estructura, que seria la que se pasaria al LLM, de forma que evitariamos
alucinaciones o errores del LLM en la lectura de la API. Esta parte esta mejor explicada en la subseccion~\ref{subsec:proteccion-anti-alucionaciones}


\subsection{Conexion con LLM}
\label{subsec:conexion-llm}

Para poder realizar la conexion con el LLM, el primer paso fue crear el módulo \textbf{utils/llm}, y dentro creamos el archivo \texttt{client.py}, el cual
implementa la funcion \texttt{chat\_completions}. Esta funcion se encarga de hacer una unica cosa: enviar un mensaje al LLM y devolver la respuesta en texto
plano.

Una decisión importante fue utilizar el estandar de la mayoria de proveedores para realizar la conexion, utilizando de esta forma el endpoint \textbf{/chat/completions}. En este caso creado para poder realizar pruebas con varios LLM de diferentes proveedores (OpenRouter, Groq, NVIDIA) y cualquier agente local.
Otra decision fue crear \texttt{.env}, de forma que la URL base, la API key y el modelo se leen de variables de entorno definidas en dicho fichero, lo que permite cambiar de proveedor sin tocar una sola línea de código.

La función mencionada anteriormente acepta tres parámetros principales:

\begin{itemize}
    \item \textbf{Mensaje de usuario}: la tarea concreta que se le envía al modelo 
    en cada llamada, en este caso el prompt con los datos del dataset, las claves 
    disponibles y los ejemplos de items reales.
    
    \item \textbf{Mensaje de sistema}: instrucciones previas que definen cómo debe 
    comportarse el modelo antes de procesar el mensaje del usuario. Se utiliza para 
    indicarle el rol que debe adoptar, el formato de respuesta esperado y lo que 
    está prohibido hacer. Es opcional, ya que técnicamente todo puede ir en el 
    mensaje de usuario, pero separarlo es una buena práctica que facilita reutilizar 
    el mismo comportamiento base en distintas llamadas.
    
    \item \textbf{Temperatura}: parámetro que controla el grado de aleatoriedad 
    en las respuestas del modelo. Un valor cercano a 0 produce respuestas más 
    deterministas y predecibles, mientras que valores más altos generan respuestas 
    más creativas pero menos consistentes. En este caso se fija por defecto en 
    \texttt{0.2}, favoreciendo la estabilidad especialmente cuando se espera JSON 
    estructurado como salida.
\end{itemize}

El primer proveedor utilizado fue OpenRouter, con el modelo de meta-llama: \\
\texttt{llama-3.3-70b-instruct:free}, de acceso gratuito y potencia más que suficiente para las primeras pruebas.


\subsection{Planner schema}
\label{subsec:planner-schema}

Otro fichero clave fue \texttt{utils/llm/planner.py} y su funcion \texttt{generate\_site\_plan}. La mision era transformar el resultado del analisis en un propuesta con campos
concretos de schema que el usuario pudiese aceptar y modificar a su gusto antes de continuar.

El schema tiene una forma fija que el resto del sistema espera:
\begin{itemize}
    \item \textbf{site\_type}: tipo de sitio recomendado, entre \texttt{blog}, \texttt{portfolio}, \texttt{catalog}, \texttt{dashboard} u \texttt{other}.
    \item \textbf{site\_title}: nombre descriptivo del sitio propuesto por el modelo.
    \item \textbf{fields}: lista ordenada de campos del dataset que el sitio debe mostrar, cada uno con su clave original y una etiqueta legible para el usuario.
\end{itemize}

Para obtener este esquema, se construye un prompt estructurado que se envía al LLM. Este prompt es el resultado de una serie de pruebas con distintas formulaciones,  
siguiendo técnicas de \emph{prompt engineering}~\cite{openai_prompt_engineering} para conseguir respuestas consistentes y parseables. La información contenida en este 
prompt es la siguiente:

\begin{itemize}
    \item \textbf{Objetivo del usuario}: texto en lenguaje natural introducido por 
    el usuario describiendo qué tipo de sitio quiere construir.
    
    \item \textbf{Claves disponibles}: lista de campos reales extraídos del dataset 
    durante la fase de análisis. El modelo solo puede proponer campos que estén 
    en esta lista.
    
    \item \textbf{Ejemplos del dataset}: hasta tres items reales del dataset, para 
    que el modelo pueda entender qué tipo de valores contiene cada campo y proponer 
    etiquetas adecuadas.
    
    \item \textbf{Contrato JSON}: estructura exacta que debe devolver el modelo, 
    con los campos obligatorios y su formato esperado.
    
    \item \textbf{Reglas de formato}: conjunto de instrucciones explícitas que 
    indican al modelo que debe devolver únicamente un objeto JSON válido, sin 
    bloques de código Markdown, sin texto introductorio y sin ningún contenido 
    fuera del JSON.
\end{itemize}

Esta rigidez sobre el formato de salida, es intencionada, ya que tras la realizacion de varias pruebas los fallos generados por los diferentes LLM siempre eran parecidos.
Como esa respuesta va a ser usada directamente, cualquier texto adicional rompe el proceso. De todas formas, debido a las alucinaciones de algunos modelos y la ignoracion 
de las restricciones en algunas ocasiones, se aplicaba una limpieza previa que eliminaba bloques Markdown y caracteres sobrantes del parseo JSON.


\subsection{Proteccion anti-alucinaciones}
\label{subsec:proteccion-anti-alucionaciones}

Los errores mas comunes siempre sucedian por la misma causa, las alucinaciones del LLM. El error mas comun era inventar informacion que no existia en la API, y que nunca se
en el parseo, por lo que no existe en el dataset contruido. Para ello hubo que aplicar varias capas de defensa contra alucinaciones para proteger los schemas.

\begin{itemize}
    \item \textbf{Validación contra claves reales}: cada clave propuesta por el LLM se comprueba contra la lista de claves reales extraída del dataset. Si la clave no existe, se intenta un match sin distinguir mayúsculas y minúsculas. Si tampoco hay coincidencia, la clave se descarta silenciosamente.
    \item \textbf{Filtro de tokens prohibidos}: existe una lista de palabras que el modelo confunde con claves del dataset porque forman parte del propio prompt, como \texttt{available\_keys}, \texttt{fields} o \texttt{examples}. Cualquier clave que coincida con esta lista se descarta automáticamente.
    \item \textbf{Detección de valores disfrazados de claves}: se detectan los casos en los que el modelo introduce un valor concreto del dataset donde debería haber puesto el nombre de un campo, como fechas en formato ISO, URLs completas o cadenas excesivamente largas.
    \item \textbf{Fallback}: si tras aplicar todas las validaciones no queda ninguna clave válida, el sistema construye automáticamente un schema de emergencia con las primeras claves del dataset, garantizando que el flujo nunca se bloquea por un fallo del modelo.
\end{itemize}

Si el parseo falla, el sistema reintenta automáticamente una vez con una pista sobre el error cometido. Si el segundo intento también falla, se activa el fallback.


\subsection{Actualizacion de modelos}
\label{subsec:actualizacion-modelos}

La integracion del LLM trajo consigo dos cambios importantes.

El primero de ellos fue una modificacion en el modelo existente \textbf{APIRequest}, añadiendole el campo \texttt{plan\_accepted}. Siendo este un booleano encargado de guardar
si el usuario ha aceptado el schema propuesto antes de continuar.

El segundo fue la creacion de un nuevo modelo llamado \textbf{GeneratedSite}, vinculado con el modelo APIRequest mediante una relacion \texttt{OneToOne}. Este modelo se
encargaba de guardar el plan aceptado, y los templates generados para el sitio web, y un identificador unico en el campo \texttt{public\_id}. Este modelo fue muy importante para
la generacion de los proyectos que se desarrollaria en la siguiente fase.


\subsection{Integracion de n8n}
\label{subsec:integracion-n8n}

En paralelo a la integración del LLM, durante esta fase se incorporó n8n al sistema por primera vez. Es una herramienta de automatización de flujos que permite conectar servicios mediante webhooks y nodos, y que más adelante sería una pieza fundamental de la arquitectura del proyecto. En este momento su uso fue todavía inicial: instalar el entorno y crear el primer flujo de pruebas.

n8n se instaló en un contenedor Docker siguiendo la imagen oficial, quedando disponible en \texttt{localhost:5678}. Una vez levantado el entorno, se creó el primer flujo en n8n y se configuró el primer webhook, apuntando a la ruta \texttt{/webhook-test/WebBuilder-Register}. El prefijo \texttt{webhook-test} indica que el webhook estaba en modo de prueba, el modo que n8n ofrece para verificar que el flujo funciona correctamente antes de activarlo en producción.

La integración en Django se realizó en \texttt{views/auth.py}, donde se añadió la función \texttt{llamar\_webhook}.
El nodo se encargaría de notificar a un nuevo usuario en base de datos, es decir, aquellos usuarios recién registrados. Django realizaba una llamada POST a ese webhook con el nombre de usuario y el email, y n8n enviaba automáticamente un correo de bienvenida. Más adelante se completaría el flujo con el webhook de login.

La función estaba diseñada para ser completamente transparente al flujo principal: si n8n no estaba disponible o la llamada fallaba, la excepción se capturaba silenciosamente y el registro continuaba con normalidad. Esto garantizaba que un fallo en el sistema de notificaciones nunca pudiera bloquear el acceso a la aplicación.


\section{Fase 4: Generador de proyectos Django}
\label{sec:generador-proytectos}

Tras la implementacion del LLM y el sistema de schema funcionando correctamente, el siguiente paso era el mas complicado del proyecto: convertir el schema aprobado por el usuario
en un proyecto django completo, funcional y desplegable. Durante esta fase del proyecto surgió el fichero más importante hasta la fecha: \textbf{utils/generator/project\_generator}, el nucleo de WebBuilder.


\subsection{Arquitectura del generador}
\label{subsec:arquitectura-generador}

El generador se organiza como un orquestador de siete pasos secuenciales, cada uno con su responsabilidad concreta y encargado de producir uno o varios ficheros del proyecto final.
El LLM interviene en los primeros seis pasos, que se encargan de la creacion de ficheros, pero en el septimo paso ya no interviene, y se concluye la ejecución con el ensamblaje de todos los ficheros generados anteriormente.

\begin{itemize}
    \item \textbf{Paso 1 — Estructura de páginas}: el LLM decide qué páginas tendrá el sitio, sus URLs, sus nombres de vista y si cada página es un listado, una página de detalle o una portada. El resultado es una lista de páginas que guiará todos los pasos siguientes.
    \item \textbf{Paso 2 — models.py}: el LLM genera el modelo de datos Django adaptado a los campos del dataset.    
    \item \textbf{Paso 3 — views.py y urls.py}: el LLM genera las vistas necesarias para cada página decidida en el paso 1. El fichero \texttt{urls.py} se genera en Python directamente a partir de la lista de páginas, sin intervención del modelo.
    
    \item \textbf{Paso 4 — base.html}: el LLM genera la plantilla base del sitio, con la estructura HTML común a todas las páginas: cabecera, navegación y bloque de contenido.
    
    \item \textbf{Paso 5 — Templates por página}: el LLM genera un template HTML independiente para cada página, adaptado a su tipo y a los campos del dataset.
    
    \item \textbf{Paso 6 — load\_data.py}: el LLM genera un comando de gestión de Django que, al ejecutarse, descarga los datos de la API original y los inserta en la base de datos del sitio generado. Este fichero es fundamental para que el sitio tenga contenido real desde el primer arranque.
    
    \item \textbf{Paso 7 — Archivos fijos}: se ensamblan todos los ficheros de infraestructura que no dependen del dataset ni del LLM: \texttt{settings.py}, \texttt{manage.py}, \texttt{wsgi.py}, \texttt{requirements.txt}, \texttt{Dockerfile} y \texttt{entrypoint.sh}.
\end{itemize}

El resultado final es un diccionario Python donde cada clave es una ruta relativa de fichero y cada valor es su contenido. Este diccionario se guarda en el campo \texttt{project\_files} del modelo \texttt{GeneratedSite}, comentado mas adelante en la subseccion~\ref{subsec:modelo-generatedsite}



\subsection{Fallbacks generacion}
\label{subsec:fallbacks-generacion}

Durante las pruebas de generacion aprendí que las alucinaciones del LLM eran constantes y que iba a llevar mucho trabajo pulir cada prompt de generacion. El LLM cometia fallos en 
cualquiera de los pasos, devolviendo codigo vacio, con errores de sintaxis o simplemente codigo random no adaptado a lo que tocaba en ese paso. Para poder tener un proyecto funcional aunque con algun error, decidí añadir \textbf{fallbacks} por cada template generado.

La funcion de los fallbacks era sencilla, si el proyecto tenia un fallo grave, el fallback actuaba y generaba una version minima pero funcional del fichero en cuestión, de esta forma el proyecto saldría funcional sin fallos graves, existirían fallos mínimos que harian que el LLM produjese un proyecto algo más genérico pero siempre desplegable.


\subsection{Evolución del modelo de datos}
\label{subsec:modelo-generatedsite}

El modelo \texttt{GeneratedSite} fue variando durante esta fase para reflejar los nuevos estados del sistema.
Se añadieron los campos \texttt{project\_files} para almacenar todos los ficheros generados, \texttt{project\_name} con el nombre del proyecto, \texttt{preview\_url} para guardar la URL del sitio desplegado (explicado en la siguiente subseccion~\ref{subsec:workflow-despliegue}), y dos nuevos conjuntos de estados. El primero, \texttt{generation\_status}, seguía el ciclo de vida de la generación con los valores \texttt{pending}, \texttt{generating}, \texttt{ready} y \texttt{error}. El segundo, \texttt{deploy\_status}, reflejaba el estado del despliegue Docker con los valores \texttt{idle}, \texttt{deploying}, \texttt{done} y \texttt{error}. 

Esta separación entre el estado de generación y el estado de despliegue permitía saber con precisión en qué punto del proceso se encontraba cada sitio.


\subsection*{Dockerfile, entrypoint y requirements}
\label{subsec:dockerfile}

Cada proyecto generado incluía un \texttt{Dockerfile}, un script \texttt{entrypoint.sh} y un fichero \texttt{requirements.txt} que en conjunto automatizaban completamente el arranque del sitio. El objetivo era que cualquier proyecto generado por WebBuilder se desplegase mediante Docker pulsando un botón, sin configuración manual por parte del usuario.

El \texttt{Dockerfile} partía de la imagen oficial \texttt{python:3.12-slim}. Copiaba el \\ \texttt{requirements.txt}, instalaba las dependencias con \texttt{pip} y copiaba el resto del proyecto. Hasta entonces el puerto expuesto era el 8000, sobre el que se serviría la aplicación, esto variaría más adelante con la mejora del flujo de despliegue de n8n.

El fichero \texttt{requirements.txt} incluido en cada proyecto generado contenía únicamente las dependencias mínimas para el funcionamiento del sitio: \texttt{django>=4.2,<6.0} como framework web, \texttt{requests>=2.28} para que el comando \texttt{load\_data} pudiera descargar los datos de la API original, y \texttt{gunicorn>=21.0} como servidor de producción. La decisión de mantener las dependencias al mínimo fue intencionada: cuantas menos dependencias tenga el proyecto generado, menos probabilidades hay de que el \texttt{docker build} falle por incompatibilidades o paquetes no disponibles.

El script \texttt{entrypoint.sh} era el responsable del arranque completo del sitio, los pasos que ejecutaba son los siguientes:

\begin{enumerate}
    \item \textbf{Generación de migraciones}: ejecutaba \texttt{python manage.py makemigrations siteapp} para crear los ficheros de migración a partir del \texttt{models.py} generado por el LLM, siendo siteapp el nombre de la aplicación interna del proyecto generado.
    \item \textbf{Aplicación de migraciones}: ejecutaba \texttt{python manage.py migrate} para crear las tablas en la base de datos SQLite del contenedor.
    \item \textbf{Carga de datos}: ejecutaba \texttt{python manage.py load\_data}, el comando generado por el LLM en el paso 6 del generador, que descargaba los datos de la API original e insertaba los registros en la base de datos.
    \item \textbf{Arranque del servidor}: lanzaba Gunicorn con dos workers y un timeout de 120 segundos para servir la aplicación Django en el puerto 8000.
\end{enumerate}

Esta secuencia garantizaba que al acceder al sitio por primera vez ya estuviera completamente funcional: base de datos creada, datos cargados y servidor en marcha. El uso de 
Gunicorn en lugar del servidor de desarrollo de Django era una decisión deliberada para que los sitios generados tuvieran un nivel mínimo de robustez en producción.


\subsection{Workflow de despligue n8n}
\label{subsec:workflow-despliegue}

En paralelo al desarrollo del generador de archivos, se fue construyendo el workflow mas importante del proyecto, encargado del despliegue de las paginas.
La decision de realizar el despliegue de las paginas con n8n y no con Django fue una decision mas o menos sencilla, el unico problema de desplegar con Django era que el proceso
de build del Docker puede tardar unos minutos, y si se ejecutaba desde Django la web, el servidor se quedaria bloqueado durante este tiempo. Al separarlo y crearlo con n8n, Django lo unico que tiene que hacer es el enviar la señal para que comience la ejecución del workflow.

Esta señal se envia desde la ruta \texttt{/webhook/webbuilder-deploy}. Django realiza una llamada POST con el nombre del proyecto como parametro y a partir de aqui el workflow se ejecuta siguiendo los siguientes nodos sucesivamente:

\begin{enumerate}
    \item \textbf{Descomprimir ZIP}: el primer nodo crea el directorio de despliegue en \\ \texttt{/files/deploys/<project\_name>} y descomprime el ZIP del proyecto generado en un subdirectorio \texttt{src/}.
    \item \textbf{Restablecer contenedor}: antes de construir la nueva imagen, el workflow intenta detener y eliminar cualquier contenedor anterior con el mismo nombre (\texttt{wb\_<project\_name>}). Si el contenedor no existe, el error se ignora y el flujo continúa sin interrupciones.
    
    \item \textbf{Elegir puerto}: el nodo escanea los puertos actualmente en uso por Docker en el rango 8100--8999 y selecciona el primer puerto libre disponible. Si no hay ningún puerto libre en ese rango, el flujo lanza un error controlado.
    
    \item \textbf{Docker Build}: ejecuta \texttt{docker build} sobre el directorio del proyecto con un timeout de 5 minutos. Si el build falla, captura los últimos 500 caracteres de los logs de error para facilitar el diagnóstico.
    
    \item \textbf{Docker Run}: arranca el contenedor en modo \emph{detached} con el nombre y el puerto asignados, mapeando el puerto elegido al puerto interno 8000 del contenedor.
    
    \item \textbf{Verificar Docker}: espera 20 segundos para dar tiempo al \texttt{entrypoint.sh} a ejecutar las migraciones y cargar los datos, y luego comprueba que el contenedor sigue en funcionamiento. Si el contenedor ha caído, recupera los últimos 40 líneas de sus logs y lanza un error con esa información para facilitar el diagnóstico.

    \item \textbf{Responder al webhook}: si todo ha ido bien, devuelve a Django un JSON con la URL del preview, el nombre del contenedor y el puerto asignado, que WebBuilder almacena en el campo \texttt{preview\_url} del modelo \texttt{GeneratedSite}.
\end{enumerate}

Este diseño garantiza que el usuario siempre recibe una respuesta clara: si el despliegue tiene éxito obtiene la URL del sitio generado, y si falla en cualquier paso Django es notificado con información suficiente para actualizar el estado del sitio a \texttt{error} y mostrar un mensaje al usuario explicando lo ocurrido.


\section{Fase 5: Refactor del proyecto y consolidación del generador}
\label{sec:refactor-consolidacion-generador}

Una vez realizadas las primeras pruebas y añadidas muchas trazas para comprobar los diferentes errores surgidos en el proceso de generacion, el codigo del 
fichero \textbf{project\_generator.py} habia crecido demasiado y se habia convertido en un fichero dificil de depurar, y aprovechando este momento de refactorizacion se realizaron mas cambios sobre gran parte de codigos del proyecto.

Esta fase se divido en dos partes importantes. En primer lugar a la refactorizacion de partes importantes, como el fichero comentado anteriormente de mucho ficheros html y css, haciendo su uso y modificacion futura mucho mas sencilla al tener todo divido en modulos.

\subsection{Refactor del generador en modulos}
\label{subsec:refactor-generador-modulos}

Estos cambios fueron muy importantes, los pasos del generador se realizaban con todo el contenido dentro del mismo fichero, algo que no era de facil escalabilidad a futuro, por lo que el primer cambio fue dividir el \textbf{project\_generator.py} en cuatro ficheros encargados cada uno de una labor
especifica dentro de \textbf{utils/generator}. Estos ficheros fueron:

\begin{itemize}
    \item \textbf{\texttt{llm\_wrappers.py}}: Centralizó todas las llamadas al LLM en tres funciones reutilizables. \texttt{llm\_call} para llamadas básicas que devuelven texto, \texttt{llm\_json\_call} para llamadas que esperan JSON como respuesta, y \texttt{llm\_call\_logged} para llamadas que además guardan un registro en base de datos para mejorar la depuracion. Esta ultima funcion fue muy importante en el proceso de generacion y depuracion, anteriormente si algo fallaba no habia forma de saber que habia pasado, pero de esta forma cada vez que se llamaba al LLM durante la generación de un proyecto, además de hacer la llamada normal, guardaba automáticamente en base de datos un registro con toda la información de esa llamada: el paso que se estaba generando, el prompt enviado, la respuesta recibida y si había habido algún error. Ese registro se guarda en \texttt{GenerationLog}
    
    \item \textbf{\texttt{fallbacks.py}}: Agrupó todos los fallbacks de emergencia en un único fichero. Cada función de este módulo generaba una versión mínima pero funcional de un artefacto concreto del proyecto: páginas, modelos, vistas o templates. Al tenerlos centralizados resultaba mucho más sencillo mejorarlos o ajustarlos sin tocar el orquestador principal.
    
    \item \textbf{\texttt{static\_files.py}}: Extrajo toda la generación de ficheros de infraestructura que no dependen del LLM: \texttt{manage.py}, \texttt{settings.py}, \texttt{wsgi.py}, \texttt{urls.py}, \texttt{Dockerfile}, \texttt{entrypoint.sh} y \texttt{requirements.txt}. El orquestador principal quedó así mucho más limpio, delegando en este módulo la construcción del scaffolding fijo.
    
    \item \textbf{\texttt{migrations\_generator.py}}: Aisló la lógica de generación automática de la migración inicial, que parseaba el \texttt{models.py} producido por el LLM y construía el fichero \texttt{0001\_initial.py} correspondiente.
\end{itemize}


\subsection{Refactor de templates y ficheros de estilos}
\label{subsec:refactor-templates-estilos}

La misma causa que creo la decision del refactor del generador ocurrio con los ficheros de presentacion del proyecto (HTML, CSS y JavaScript).
Hasta el momento los templates HTML eran ficheros que incluian todas las secciones en un unico bloque, al igual que ocurria con los ficheros CSS y JavaScript. Como comentaba anteriormente, esta fase introdujo una reorganizacion completa de la interfaz con la separacion de responsabilidades.

En cuanto a los templates se creo un nuevo modulo \textbf{partials}, encargado de contener los fragmentos que representan una seccion concreta de cada pagina. Cada fichero principal paso a contener sus partials, incluyendose mediante la etiqueta \textbf{include} de Django.

La home del proyecto, que hasta entonces era el fichero mas complejo, se dividió en seis partials independientes dentro de \texttt{templates/WebBuilder/partials/home}, quedando asi la nueva estructura:

\begin{itemize}
    \item \texttt{hero.html} — la sección principal de bienvenida
    \item \texttt{examples.html} — ejemplos de sitios generados
    \item \texttt{how.html} — explicación del funcionamiento
    \item \texttt{llms.html} — sección de modelos de lenguaje compatibles
    \item \texttt{metrics.html} — métricas públicas del sistema
    \item \texttt{cta.html} — botón de acceso a la aplicación
\end{itemize}

El asistente se dividió en tres partials dentro de \texttt{templates/WebBuilder/partials/assistant/}:

\begin{itemize}
    \item \texttt{input.html} — formulario del paso 1, donde el usuario 
    introduce la URL de la API o sube un fichero local para su análisis
    \item \texttt{analysis.html} — resultado del paso 2, que muestra las 
    estadísticas del dataset analizado: número de items encontrados, campos 
    disponibles y ruta de la colección principal
    \item \texttt{preview.html} — resultado del paso 3, donde se muestra 
    el schema propuesto por la IA y el usuario puede revisarlo y aprobarlo 
    antes de continuar
\end{itemize}

El historial se dividió en dos partials dentro de \texttt{templates/WebBuilder/partials/history/}:

\begin{itemize}
    \item \texttt{published\_sites.html} — listado de sitios ya generados 
    y desplegados por el usuario
    \item \texttt{recent\_analysis.html} — listado de análisis recientes 
    realizados por el usuario
\end{itemize}

Los ficheros estáticos sufrieron una reorganización equivalente. 
El CSS pasó a estar organizado por página y por sección dentro de 
\texttt{static/css/}:

\begin{itemize}
    \item \texttt{home/hero.css} — estilos de la sección principal de bienvenida
    \item \texttt{home/examples.css} — estilos de la sección de ejemplos
    \item \texttt{home/how.css} — estilos de la sección de funcionamiento
    \item \texttt{home/llms.css} — estilos de la sección de modelos compatibles
    \item \texttt{home/metrics.css} — estilos de la sección de métricas públicas
    \item \texttt{home/cta.css} — estilos de la sección de acceso a la aplicación
    \item \texttt{assistant.css} — estilos del asistente de generación
    \item \texttt{history.css} — estilos del historial del usuario
    \item \texttt{edit.css} — estilos de la página de edición del schema
    \item \texttt{login.css} — estilos de las páginas de login y registro
    \item \texttt{navbar.css} — estilos de la barra de navegación
    \item \texttt{backgrounds.css} — fondos y efectos visuales compartidos
\end{itemize}

El JavaScript siguió el mismo criterio con ficheros independientes 
por página en \texttt{static/js/}:

\begin{itemize}
    \item \texttt{home.js} — lógica interactiva de la home
    \item \texttt{assistant.js} — lógica del flujo del asistente
    \item \texttt{edit.js} — lógica de la edición del schema
    \item \texttt{site\_render.js} — lógica de la vista previa del sitio generado
\end{itemize}

El \texttt{base.html} se actualizó para cargar las fuentes de Google Fonts 
necesarias para el nuevo diseño: Inter para el texto general, 
y Geist y Geist Mono para elementos de código y datos técnicos.

Esta reorganización, aunque invisible para el usuario final, facilitó enormemente el mantenimiento y la evolución de la interfaz en los sprints siguientes, ya que cualquier cambio en una sección concreta de una página solo requería tocar un fichero pequeño y bien delimitado.


\subsection{Enriquecimiento de prompts}
\label{subsec:enriquecimiento-prompts}

Tras unas cuantas pruebas, llegue a la conclusion de un problema importante, podria ocurrir que un usuario no familiarizado con los prompts de IA no supiese como hacerlo adecuadamente, o que incluso lo dejase vacio. Para resolver este problema se creo el fichero \texttt{utils/llm/enrich\_prompts.py}, 
encargado de analizar automaticamente los campos del dataset antes de construir el prompt y añadir contenido personalizado.
Si el dataset tenia campos con imagenes se le indica al modelo que las usara de forma destacada. Si tenia precios, que los resalte visualmente. Si tenia
fechas que las mostrara de forma clara. De esta forma se garantiza que el LLM siempre recibiera suficiente contexto para tomar decisiones razonadas.
De todos modos si algo no estuviese en la mente del usuario, en el siguiente paso de editar el schema podria modificarlo añadiendo los campos necesarios o modificando los existentes.


\subsection{Trazabilidad: GenerationLog}
\label{subsec:trazabilidad-generationlog}

Anteriormente en las generaciones si algo fallaba no había forma de saber qué prompt se había enviado exactamente, qué había respondido el modelo ni si había habido reintentos. Para arreglar esto se creo el modelo \textbf{GenerationLog}, vinculado mediante una clave a su correspondiente \textbf{GeneratedSite}. 

Cada llamada al LLM durante la generación de un proyecto producía un registro que almacenaba el paso concreto que se estaba generando, el modelo de LLM utilizado, el prompt de sistema y el prompt de usuario enviados, la respuesta en crudo devuelta por el modelo, los errores de consistencia detectados y si la llamada había necesitado un reintento. Esto permitía analizar a posteriori qué pasos fallaban con más frecuencia, qué modelos producían mejores resultados y qué tipos de datasets generaban más inconsistencias.


\subsection{Verificacion de consistencia}
\label{subsec:verificacion-consistencia}

Otro problema frecuente era que el LLM no tenia suficiente contexto entre ficheros, principal causa de las alucinaciones del LLM. Ocurria lo siguiente:
se definia un campo con un nombre en \textbf{models.py} y luego en los demas ficheros se referenciaba con ligeros cambios, provocando errores de ejecución.
Para detectar estos casos antes de que el usuario desplegase el proyecto, se implemento el fichero \texttt{utils/llm/consistency\_checker.py}.

Este fichero extraía mediante expresiones regulares los campos definidos en \texttt{models.py} y buscaba en los demas templates las referencias a esos campos. Cualquier referencia encontrada diferente que no correspondiera al campo real del modelo, se registraba como error de consistencia. 
Estos errores se guardaban en el campo \texttt{consistency\_errors} del modelo \texttt{GenerationLog} correspondiente, permitiendo identificar que pasos producian mas inconsistencias y como se producian, para el consiguiente arreglo a futuro.


\subsection{Panel de métricas}
\label{subsec:panel-metricas}

Para tener un uso mas sencillo y accesible del contenido de GenerationLog añadi una vista de metricas en \textbf{views/metrics.py}, accesible unicamente para usuarios con permisos de administrador. El objetivo de esta vista era tener visibilidad real sobre el comportamiento del generador, los fallos mas recurrentes, que modelos eran mas optimos para cada tipo de pagina, etc.

El panel mostraba las siguientes estadísticas agregadas:

\begin{itemize}
    \item \textbf{Resumen general}: número total de registros en \texttt{GenerationLog}, sitios generados y análisis realizados hasta la fecha.
    \item \textbf{Desglose por paso}: para cada uno de los siete pasos del generador, se mostraba el número de llamadas realizadas, cuántas habían necesitado reintento y la tasa de reintento resultante. Esto permitía identificar qué pasos eran más problemáticos y priorizar las mejoras.
    \item \textbf{Distribución por modelo}: agrupación de las llamadas por modelo de LLM utilizado, lo que facilitaba comparar el rendimiento de distintos proveedores cuando se probaban configuraciones diferentes.
\end{itemize}

Esta informacion resulto util para mejorar las siguientes fases donde cambios minimos producirian grandes resultados.


\section{Fase 6: Funcionalidades avanzadas}
\label{sec:funcionalidades-avanzadas}

Una vez concluida una buena base del sistema de generacion y despliegue, esta fase se dedico a implementar funcionalidades que ampliaban las capacidades del sistema en tres direcciones: historial, ingestion de datos y LLM personalizado

En primer lugar, se añadieron nuevos modelos para cubrir las necesidades que tenia el historial, se crearon \texttt{SiteVersion}, para mantener 
un historial completo de todas las versiones de ficheros de cada sitio generado y poder recuperar cualquier estado anterior, y \texttt{SiteUser}, 
para gestionar los usuarios internos de cada proyecto Django generado directamente desde WebBuilder, pudiendo asi introducir usuarios o superusuarios.

En segundo lugar, el módulo de ingestión se amplió para soportar dos nuevos formatos de entrada muy habituales en APIs de datos abiertos: CSV y GeoJSON. Detectados de forma automatica como ocurria con XML y JSON. Normalizando los datos a una lista de diccionarios (\texttt{list[dict]}) como ocurria con los anteriores tipos, para su posterior uso.


Por último, se implementó el modelo \texttt{UserProfile}, que permite a cada usuario configurar su LLM local, con su propia API key cifrada, eliminando la dependencia de seleccion cerrada de LLMs y abriendo el sistema a cualquier modelo compatible con el estándar OpenAI.

\subsection{Nuevos formatos de entrada: CSV y GeoJSON}
\label{subsec:nuevos-formatos-entrada}

Hasta este momento el modulo de ingestion solo soportaba JSON y XML. En esta fase se ampliaron los formatos añadiendo soporte para CSV y GeoJSON, que son dos de los formatos mas habituales en APIs publicas.

La deteccion del formato se añadio en el archivo existente \textbf{utils/ingest/parsers.py}. 

El formato GeoJSON se detectaba antes que el formato JSON, ya que es un formato mas restrictivo. Se realiza inspeccionando el campo \texttt{type} del objeto raíz: si era \texttt{FeatureCollection}, el fichero se procesaba como GeoJSON. Mientras que para el formato CSV se inspeccionaba la primera linea en busca de separadores de coma, o punto y coma, habituales en este tipo de ficheros.

El parseo de cada formato tenía su propia función:

\begin{itemize}
    \item \textbf{CSV}: se utilizó la librería estándar \texttt{csv} de Python con detección automática del dialecto mediante \texttt{csv.Sniffer}. La lectura se limitaba a 2000 filas para evitar problemas de memoria con ficheros muy grandes.
    \item \textbf{GeoJSON}: cada \texttt{Feature} del \texttt{FeatureCollection} se aplanaba extrayendo sus propiedades como un diccionario plano, añadiendo además el tipo de geometría como campo \texttt{geometry\_type}. De esta forma el resto del sistema podía tratar el resultado igual que cualquier otro JSON.
\end{itemize}

\subsection{Ejemplos few-shot por tipo de sitio}
\label{subsec:ejemplos-fewshot}

Otro problema detectado en las generaciones era que los templates HTML producidos por el LLM tendían a ser genéricos y con poco estilo visual. Para mejorar la calidad del código generado se introdujo la carpeta \texttt{utils/llm/examples/}, con ejemplos de templates HTML reales por tipo de sitio: \texttt{blog.py}, \texttt{catalog.py}, \texttt{dashboard.py} y \texttt{portfolio.py}. Estos ejemplos se incluían en el prompt del paso 5 del generador como referencia de estilo, siguiendo la técnica de \emph{few-shot prompting}~\cite{openai_prompt_engineering}, que consiste en proporcionar al modelo ejemplos concretos del resultado esperado para que los imite en lugar de inventar una solución desde cero. El resultado fue una mejora notable en la calidad visual y estructural de los templates generados.


\subsection{Versionado de proyectos generados}
\label{subsec:versionado-paginas}

Hasta este momento, generar o editar una API existente en el historial, causaba que los ficheros existentes se sobreescribieran, de esta forma era imposible tener un control de versiones o recurrir a una prueba anterior que fuese mejor. Para solucionar esto se creo el modelo \textbf{SiteVersion}, enlazado mediante una clave foranea a \textbf{GeneratedSite}. Cada version estaba nombrada con un numero incremental, y almacenaba una copia completa de los ficheros, una etiqueta descriptiva opcional y la fecha de creacion del proyecto. Las versiones estaban ordenadas de forma mas reciente a mas antigua, permitiendo asi al usuario consultar un historial funcional y la posibilidad de recuperar cualquier version anterior.

\subsection{Usuarios del proyecto}
\label{subsec:usuarios-proyecto}

Cada sitio Django generado incluia su propio sistema de autenticacion, son paginas de login y register. Para gestionar los usuarios de esos sitios desde la propia plataforma se creo el modelo \textbf{SiteUser}, vinculado a \textbf{GeneratedSite}. Cada usuario del sitio tenia su nombre, una contraseña y un rol. Este ultimo campo podria ser \texttt{super} para la creacion de un superusuario con acceso al panel de administracion Django del proyecto creado, o \texttt{normal} para un usuario con acceso solo al frontend. La combinacion de proyecto y nombre de usuario era unica, evitando asi duplicados dentro del mismo proyecto.

\subsection{LLM local y catalogo de modelos disponibles}
\label{subsec:llm-local-catalogo-modelos}

Hasta este momento el modelo de LLM utilizado para generar páginas era fijo, configurado en el fichero \textbf{.env} del servidor, sin que el usuario pudiera cambiarlo. En esta fase se implementaron dos cambios importantes respecto a esto.

En primer lugar, se creó el fichero \texttt{utils/llm/llm\_catalog.py}, que definía un catálogo de modelos conocidos y recomendados. Cada entrada del catálogo incluía el identificador del modelo, su nombre legible, el proveedor, la URL base del API, una descripción y una lista de ventajas e inconvenientes. En este sprint el catálogo incluía tres modelos:

\begin{itemize}
    \item \textbf{Llama 4 Scout} (Groq) — modelo rápido y eficiente, recomendado para generaciones ágiles con buena calidad de código
    \item \textbf{Llama 3.3 70B} (OpenRouter) — modelo más potente, recomendado para webs complejas que requieren mayor capacidad de razonamiento
    \item \textbf{Qwen3 Coder} (OpenRouter) — modelo especializado en código, especialmente bueno generando HTML y CSS limpio
\end{itemize}


En segundo lugar, se implementó un selector de modelo completo, accesible desde el inicio del proceso, justo antes de analizar la URL. El selector ofrecía tres modos de uso:

\begin{itemize}
    \item \textbf{Predeterminado}: el sistema selecciona automáticamente el 
    modelo más adecuado según el dataset analizado, buscando el equilibrio 
    óptimo entre velocidad y calidad. Es la opción recomendada para la 
    mayoría de casos.
    
    \item \textbf{Elegir modelo}: despliega el catálogo de modelos disponibles 
    definido en \\ \texttt{utils/llm/llm\_catalog.py}, mostrando para cada modelo 
    su nombre, proveedor, descripción y sus ventajas e inconvenientes.
    
    \item \textbf{Modelo personalizado}: utiliza el LLM configurado por el 
    usuario en su perfil. Si el usuario no tiene ninguno configurado, esta 
    opción aparece deshabilitada.
\end{itemize}

El catálogo cubría los casos más habituales, pero existía la necesidad de 
permitir que cualquier usuario pudiera conectar su propio proveedor de LLM, 
ya fuera uno privado, uno de pago con mejores capacidades o simplemente uno 
diferente a los disponibles en el catálogo. Para soportar esto se creó el 
modelo \texttt{UserProfile}, vinculado en relación \texttt{OneToOne} con el 
usuario de Django. El modelo almacenaba tres campos opcionales:

\begin{itemize}
    \item \textbf{URL base del proveedor} (\texttt{custom\_llm\_base\_url}): 
    dirección del endpoint del proveedor OpenAI-compatible, por ejemplo 
    \texttt{https://api.groq.com/openai/v1} para Groq o 
    \texttt{https://openrouter.ai/api/v1} para OpenRouter.
    
    \item \textbf{Identificador del modelo} (\texttt{custom\_llm\_model}): 
    nombre exacto del modelo tal y como lo espera el proveedor, por ejemplo 
    \texttt{meta-llama/llama-3.3-70b-instruct:free}.
    
    \item \textbf{API key} (\texttt{custom\_llm\_api\_key}): clave de acceso 
    al proveedor, cifrada en base de datos mediante \texttt{EncryptedCharField} 
    para que nunca se almacene en texto plano.
\end{itemize}


\section{Fase 7: Pulido y cierre}
\label{sec:pulido-cierre}

Esta ultima fase se dedico a pulir el sistema en todos sus puntos débiles, que eran: mejorar la calidad visual de los sitios generados, consolidar la infraestructura de produccion y completar los workflows de n8n que quedaban pendientes. Con infraestructura de produccion me refiero a los cambios que hacen que el proyecto sea util fuera del entorno de desarrollo, como por ejemplo la migración a una base de datos más robusta como PostgreSQL para soportar accesos concurrentes de multiples usuarios, la gestión automática de contenedores Docker activos o la monitorización diaria del sistema. Es decir, no se añadieron funcionalidades realmente nuevas, sino que se llevaron las existentes a un nivel de madurez adecuado para la entrega.


\subsection{Sistema de presets visuales}
\label{subsec:sistema-presets-visuales}

Uno de los principales problema durante todo el desarrollo fue que los proyectos tendian a parecerse entre si visualmente, independiente del tipo de contenido de la API, del prompt del usuario y del sitio generado. Para resolverlo se creó el módulo \textbf{utils/llm/design}, organizado en tres ficheros con responsabilidades distintas. Estos ficheros eran presets.py, snippets.py y theme\_rules.py.

El fichero \texttt{presets.py} define un catálogo de estilos visuales completos y muy distintos entre sí. Cada preset describe de forma exhaustiva la paleta de colores, el tipo de layout, la estética de los componentes y las palabras clave del prompt del usuario que deben activarlo. Los presets disponibles en esta fase son:

\begin{itemize}
    \item \textbf{Neo Terminal}: estética cyberpunk con fondo casi negro y acento neón ácido. Layout denso orientado a escaneo rápido de información. Activado por palabras como \emph{terminal}, \emph{dark}, \emph{gaming} o \emph{futurista}.
    \item \textbf{Bento Pop}: rejilla asimétrica con bloques de colores vivos e ilustrativos. Activado por palabras como \emph{bento}, \emph{mosaico} o \emph{asimétrico}.
    \item \textbf{Quiet Editorial}: tipografía protagonista, márgenes amplios y estética de revista cultural sobria. Preset por defecto para sitios de tipo blog.
    \item \textbf{Magazine Split}: portada con gran imagen dividida y tipografía impactante. Preset por defecto para catálogos.
    \item \textbf{Glass Orbit}: estética glassmorphism con fondos oscuros y elementos traslúcidos. Preset por defecto para portfolios.
\end{itemize} 

Cada tipo de sitio tiene asignado un preset por defecto que se aplica cuando el prompt del usuario no contiene ningun estilo descrito, garantizando que el resultado siempre tenga una identidad visual coherente con el tipo de contenido.

El fichero \texttt{snippets.py} complementa los presets con fragmentos de HTML concretos y válidos para cada estilo: cards, heros y filas de listado con las clases Tailwind exactas que corresponden a cada preset. Estos snippets se incluyen en los prompts de generación de templates como ejemplos de referencia visual, asegurando que el LLM produzca código coherente con el estilo elegido.

El fichero \texttt{theme\_rules.py} centraliza tokens visuales reutilizables como clases base para contenedores, cards, botones, badges y tipografía, junto con reglas críticas de uso de color para evitar que el LLM genere clases Tailwind inválidas, como \texttt{bg-\#111111} en lugar de la sintaxis correcta \texttt{bg-[\#111111]}.

\subsection{Migración a PostgreSQL}
\label{subsec:migracion-postgresql}

Durante todo el desarrollo, WebBuilder había utilizado SQLite como base de datos, lo cual era suficiente para el entorno de desarrollo pero inadecuado para producción. En este sprint se realizó la migración a PostgreSQL, que ofrece mayor robustez, soporte para accesos concurrentes y mejor rendimiento con volúmenes de datos crecientes.

La configuración en \texttt{settings.py} se actualizó para leer los parámetros de conexión desde variables de entorno: nombre de la base de datos (\texttt{DB\_NAME}), usuario (\texttt{DB\_USER}), contraseña (\texttt{DB\_PASSWORD}), host (\texttt{DB\_HOST}) y puerto (\texttt{DB\_PORT}, por defecto \texttt{5432}). La configuración anterior de SQLite se migró a la nueva base de datos, teniendo los logs y métricas anteriores.


\subsection{Workflow generation-done}
\label{subsec:workflow-generation-done}

Con la base de notificaciones ya preparada desde el fichero \\ \texttt{utils/generator/notifications.py}, se implementó el workflow de n8n \emph{generation-done}. Este workflow se activa mediante un webhook al que Django llama en cuanto termina la generación de un sitio, enviando un payload con los datos relevantes: email y nombre del usuario, título y tipo del sitio generado, URL de la API original, URL del preview si ya está desplegado, URL de descarga del ZIP, número de campos e items procesados, duración de la generación en segundos y modelo de LLM utilizado.
n8n recibe este payload y envía automáticamente un correo de notificación al usuario informándole de que su sitio está listo. La llamada al webhook se realizó de forma completamente transparente al flujo principal: si n8n no estaba disponible o la llamada fallaba, la excepción se capturaba silenciosamente mediante un bloque \texttt{try/except} y la generación continuaba con normalidad.


\subsection{Workflow health-check}
\label{subsec:workflow-health-check}

El workflow de \emph{health-check} se configuró para ejecutarse automáticamente cada día mediante un nodo \emph{Schedule Trigger} de n8n. Su función es verificar que el sistema sigue operativo realizando una petición HTTP a la URL principal de WebBuilder y comprobando que la respuesta es correcta. Si la verificación falla, n8n envía una alerta por correo al administrador del sistema con información sobre el error detectado. Este workflow fue especialmente útil durante las últimas fases de desarrollo para detectar caídas del servidor sin necesidad de monitorización manual.


\subsection{Workflow auto-shutdown}
\label{subsec:workflow-auto-shutdown}

El workflow de \emph{auto-shutdown} resolvió un problema operativo importante: los contenedores Docker de los sitios generados permanecían activos indefinidamente aunque nadie los estuviera usando, consumiendo recursos del servidor. Este workflow se ejecuta de forma programada cada hora, consulta la lista de contenedores activos con el prefijo \texttt{wb\_} y detiene aquellos que llevan inactivos más de un tiempo configurable. Esto permite mantener el servidor limpio y los recursos disponibles para nuevas generaciones sin intervención manual.


\subsection{Mejora del workflow de despliegue}
\label{subsec:mejora-workflow-despliegue}

El workflow de despliegue que se había construido en el Sprint 4 carecía de un mecanismo centralizado de gestión de errores. Si cualquiera de los nodos fallaba, la petición de Django quedaba sin respuesta. En este sprint se añadió el nodo \textbf{Responder Error}, que actúa como rama de error compartida para todos los nodos del workflow. Cuando cualquier paso falla, este nodo captura el error y devuelve a Django una respuesta JSON estructurada con el código HTTP 500, un mensaje descriptivo del error y el paso en el que ocurrió. Esto permite que Django actualice el estado del sitio a \texttt{error} y muestre un mensaje útil al usuario en lugar de dejar la operación colgada indefinidamente.



% Fin diseño e implementacion
%%%%%%%%%%%%%%%%%%%%%%%%%%%%%%%%%%%%%%%%%%%%%%%%%%%%%%%%%%%%%%%%%%%%%%%%%%%%%%%%


